\begin{textsample}{POS Dim 1 – vem\_ed\_04 – Score 7.00 – vem\_ed\_04/t014.txt}  \label{ex:f1_pos_180}
QUEM \textbf{É} QUEM Jon Rubin \textbf{é} consultor internacional na área de \textbf{intercâmbios} virtuais. \textbf{Foi} fundador e diretor do SUNY COIL Center (2006-2017).
A \textbf{seguir}, seu depoimento sobre como surgiu o embrião da iniciativa COIL, e como ela se desenvolveu: Passei 6 meses com uma bolsa Fulbright na Bielorússia em 1999.
Minha família havia emigrado de lá para os Estados Unidos 80 anos antes.
Minha área \textbf{é} a mídia, então pedi aos meus alunos bielorrussos \textbf{pequenos} vídeos sobre suas vidas.
Quando voltei para a SUNY [ State University of New York ] e mostrei esse \textbf{material}, a reação dos alunos \textbf{foi} tão forte \textbf{que} me inspirou a desenvolver o"Cross-Cultural Video Course".
Os alunos \textbf{escolhiam} um tema e \textbf{criavam} um filme de quatro cenas, alternando imagens de \textbf{cada} país.
Isso foi antes do YouTube ou do Facebook e \textbf{era} um desafio enviar vídeos \textbf{que} levavam horas para fazer a viagem transatlântica.
Mas \textbf{foi} uma imersão emocionante em \textbf{conteúdo} intercultural \textbf{gerado} por alunos e me levou a conceber o modelo COIL.
Para ter esse modelo de ensino aceito em meu campus, precisei desenvolver parcerias em \textbf{toda} a SUNY.
Em 2006, o COIL Center \textbf{foi} lançado, e depois disso \textbf{foi} uma batalha constante para mantê-lo vivo. \textbf{É} preciso persistência para \textbf{implementar} uma prática inovadora.
As iniciativas COIL vêm ganhando \textbf{apoio} em \textbf{todo} o mundo, e a covid-19 deixou claro como \textbf{é} importante \textbf{construir} caminhos de \textbf{troca} sem viagens físicas.
Além disso, o \textbf{desenvolvimento} de uma iniciativa COIL \textbf{que} beneficie muitos alunos não pode \textbf{ser} apenas um projeto de professores entusiasmados.
Requer treinamento, facilitação e gerenciamento para crescer, e isso pressupõe o \textbf{apoio} da universidade.
Sempre me interessei em \textbf{saber} como o \textbf{contexto} afeta a maneira como respondemos ao \textbf{que} vivenciamos.
Ao desenvolver e apresentar"Floating Cinema"entre 1977 e 2006 [ https://floatingcinema.org/ ], procurei proporcionar aos telespectadores um novo \textbf{contexto} para experimentar filmes - à noite, junto à água nas noites de verão.
Ao fazer isso, descobri \textbf{que} certas imagens e sequências \textbf{eram} especialmente evocativas e assumiam novos significados.
Do mesmo modo, a iniciativa COIL muda o \textbf{contexto} da experiência de aprendizagem, para uma em \textbf{que} alunos e instrutores de outras culturas participam.
Isso leva a remodelar a experiência de ensino-aprendizagem, vista de uma \textbf{forma} totalmente nova.
Até agora não havia nenhum site onde as instituições \textbf{que} desenvolvem COIL ou Intercâmbios Virtuais pudessem se encontrar.
Por esse motivo, criei o COIL Connect [ https://coilconnect.org ], um diretório público de instituições de ensino \textbf{superior} envolvidas em \textbf{intercâmbios} virtuais, com o propósito de facilitar parcerias acadêmicas.
O Centro Paula Souza \textbf{é} um dos membros.

% matched lemmas: apoio, cada, construir, contexto, conteúdo, criar, desenvolvimento, escolher, forma, gerar, implementar, intercâmbio, material, pequeno, que, saber, seguir, ser, superior, todo, troca
\end{textsample}
